\documentclass[11pt]{article}
\usepackage{amsmath,amssymb}
\usepackage[margin=1in]{geometry}
\usepackage{hyperref}
\usepackage{booktabs}
\emergencystretch=3em

\title{Zero-phase dictionary for unequal exponents}
\author{Claude, with Daniel Murfet}
\date{2026-09-07}

\begin{document}
\maketitle
\sloppy

\begin{abstract}
Unit 198, the first (and, it turns out, only necessary) unit of Astra~\#20's programme~B. The
quadratic-kernel chart programme (\texttt{twoDGeneral}, parameter $n$, kernel $e^{-\beta n s^2+\beta\sqrt n\,s a}$)
and the monomial normal-moment programme (\texttt{monomialBoxReal}, kernel $e^{-\beta N u^{2k}}$) are
tied together at zero phase and \emph{unequal} exponent ratios $p_1=(h_1+1)/k_1<p_2=(h_2+1)/k_2$:
the unit-cutoff zero-phase chart integral \emph{is} the two-dimensional unit-box monomial integral at
the same parameter, so the mixed-ratio monomial theorem transports to
\[
n^{p_1/2}\int_0^1\!\!\int_0^1u^{h_1}v^{h_2}e^{-\beta n(u^{k_1}v^{k_2})^2}\,dv\,du
\longrightarrow\frac{\Gamma(p_1/2)}{2k_1\,\beta^{p_1/2}\,(h_2+1-p_1k_2)}
\]
(\texttt{headline\_quadratic\_zeroPhase\_mixed}, plus the swapped version), and uniqueness of limits
against the quadratic programme's own no-log theorem evaluates its constant in closed form:
$\mathrm{noLogConst}(\beta,0,p_1,p_2)=\Gamma(p_1/2)\beta^{-p_1/2}/(2(p_2-p_1))$. Zero
\texttt{sorry}/\texttt{axiom}; 9075 jobs.
\end{abstract}

\section{The things formalised}
{\small
\begin{verbatim}
theorem twoDGeneral_zeroPhase_eq_twoDAmp (hn : 0 ≤ n) :
    twoDGeneral β 0 b n h₁ h₂ k₁ k₂ = twoDAmp β b (√n) h₁ h₂ k₁ k₂ (fun _ _ _ => 1)
theorem twoDGeneral_zeroPhase_eq_monomial (hn : 0 ≤ n) :
    twoDGeneral β 0 1 n h₁ h₂ k₁ k₂ = monomialBoxReal 2 ![h₁, h₂] ![k₁, k₂] β n
theorem multCount_two_mixed / monomialMixedConst_two_mixed :
    monomialMixedConst ![h₁, h₂] ![k₁, k₂] l β = Γ l * β^(-l) / (2k₁ (h₂ + 1 - 2k₂ l))
theorem monomial_two_mixed_tendsto (hlt : p₁ < p₂) :
    Tendsto (fun n => monomialBoxReal 2 ![h₁,h₂] ![k₁,k₂] β n / n^(-(p₁/2))) atTop
      (𝓝 (monomialMixedConst ![h₁,h₂] ![k₁,k₂] (p₁/2) β))
theorem headline_quadratic_zeroPhase_mixed (hlt : p₁ < p₂) :
    Tendsto (fun n => n^(p₁/2) * twoDGeneral β 0 1 n h₁ h₂ k₁ k₂) atTop
      (𝓝 (Γ (p₁/2) / (2k₁ * β^(p₁/2) * (h₂ + 1 - p₁ k₂))))
theorem headline_quadratic_zeroPhase_mixed' (hgt : p₂ < p₁) : ... (coordinates exchanged)
theorem noLogConst_zeroPhase_eq_monomialMixedConst (hlt : p₁ < p₂) :
    1 / (k₁ k₂) * noLogConst β 0 p₁ p₂ = monomialMixedConst ![h₁,h₂] ![k₁,k₂] (p₁/2) β
theorem noLogConst_zeroPhase_eq (hlt : p₁ < p₂) :
    noLogConst β 0 p₁ p₂ = Γ (p₁/2) * β^(-(p₁/2)) / (2 (p₂ - p₁))
\end{verbatim}
}

\section{Mathematics}
\paragraph{Dictionary.} With $a=0$ the quadratic kernel is $e^{-\beta n(u^{k_1}v^{k_2})^2}=e^{-\beta n u^{2k_1}v^{2k_2}}$,
so the monomial parameter is $n$ itself (not $n^2$, contrast unit~193 where \texttt{twoDAmp} is
parametrised by $N=\sqrt n$). The proof goes through unit~193's
\texttt{twoDAmp\_const\_eq\_monomialCutoff} at $N=\sqrt n$ and $(\sqrt n)^2=n$.

\paragraph{Transport.} The monomial mixed theorem at $\lambda=p_1/2$ has $J=\{0\}$, $|J|=1$, so the
logarithmic power is $(\log n)^0$ and the constant is
$\Gamma(\lambda)\beta^{-\lambda}/0!\cdot\frac1{2k_1}\cdot\frac1{h_2+1-2k_2\lambda}$. No factor of two
arises anywhere (no square in the parameter, no logarithm), which is the audit Astra~\#20 asked for.

\paragraph{Constant cross-validation.} The quadratic programme proves
$\texttt{twoDGeneral}\sim\frac1{k_1k_2}b^{k_2(p_2-p_1)}\,\mathrm{noLogConst}(\beta,a,p_1,p_2)\,n^{-p_1/2}$
with $\mathrm{noLogConst}(\beta,a,p_1,p_2)=\int_0^\infty x^{p_1-p_2-1}\int_0^xt^{p_2-1}e^{-\beta t^2+\beta ta}\,dt\,dx$.
At $a=0$, $b=1$ both programmes describe the same function, so
$\frac1{k_1k_2}\mathrm{noLogConst}(\beta,0,p_1,p_2)=\frac{\Gamma(p_1/2)\beta^{-p_1/2}}{2k_1k_2(p_2-p_1)}$
(using $h_2+1=p_2k_2$), i.e.\ $\mathrm{noLogConst}(\beta,0,p_1,p_2)=\Gamma(p_1/2)\beta^{-p_1/2}/(2(p_2-p_1))$.
This is exactly what Fubini gives directly ($\int_t^\infty x^{p_1-p_2-1}dx=t^{p_1-p_2}/(p_2-p_1)$,
then $\int_0^\infty t^{p_1-1}e^{-\beta t^2}dt=\tfrac12\beta^{-p_1/2}\Gamma(p_1/2)$), but the Lean proof
is by uniqueness of limits and so is valid only for ratios realised by natural data. Numerical check
at $\beta=1.3$, $(h,k)=((0,2),(1,1))$, $p_1=1$, $p_2=3$: quadrature $0.38863621594708$ vs closed form
$0.38863621594708$; $n^{1/2}Q_n$ at $n=10^6$ is $0.388636$.

\section{Lean notes}
\begin{itemize}
\item \texttt{ring} cannot prove $\frac{X}{2k_1(h_2+1-2k_2l)}=X\cdot\frac1{2k_1}\cdot\frac1{h_2+1-2k_2l}$
  because it expands the product inside the inverse into a sum first and then treats the inverse of
  the sum as an atom. \texttt{simp only [div\_eq\_mul\_inv, mul\_inv, one\_mul]} before \texttt{ring}
  distributes the inverse over the product syntactically.
\item Rewriting a \texttt{≠} goal \texttt{(h₂+1)/(2k₂) ≠ p₁/2} with \texttt{← hp₂} fails (no $p_2$
  present); establish the equality \texttt{ratioExp … 1 = p₂/2} first and rewrite with it.
\item \texttt{monomialBoxReal\_mixed\_tendsto 1 ![h₁,h₂] ![k₁,k₂]} elaborates without fuss
  (\texttt{Fin (1+1)} vs \texttt{Fin 2}).
\end{itemize}

\section{Status}
Astra~\#20 programme~B is closed in one unit (the general mixed-$d$ transport it mentioned as a
possible follow-on is already available in the monomial programme; the quadratic side has no
arbitrary-dimensional mixed theorem beyond one dominated coordinate, and nothing in the paper
requires one at zero phase). Next: gate for programme~A (general-$d$ random-phase leading term).
\end{document}
